% ID: SEL-D0034-OB-E04
% ESTADO: cotejada. Segunda pasada visual y borrador compilado revisados.
% Fuente: Selectividad/matematicas_aplicadas/titular_septiembre_1_MATEMATICAS_APLICADAS.pdf
% SHA-256: c637bbde5c26c0b35b59202bd832ddc97990d3845365ce483fbf2a49832f73fe
% Página: 2; opción B; ejercicio original 4.

\begin{enumerate}[label=\alph*)]
\item \textbf{(1.25 puntos)} Se desea tomar una muestra aleatoria estratificada de las personas mayores de edad de un municipio, cuyos estratos son los siguientes intervalos de edades, en años: de 18 a 30, de 31 a 45, de 46 a 60 y mayores de 60. En el primer intervalo hay 7500 personas, en el segundo hay 8400, en el tercero 5700 y en el cuarto 3000. Calcule el tamaño de la muestra total y su composición, sabiendo que el muestreo se hace con afijación proporcional y se han elegido al azar 375 personas del primer estrato.
\item \textbf{(1.25 puntos)} Dada la población $\{2,4,6\}$ construya todas las muestras posibles de tamaño 2, que se puedan formar mediante muestreo aleatorio simple, y halle la varianza de las medias muestrales de todas las muestras.
\end{enumerate}
